\section{Introduction}

For a graph with nonnegative edge weights, Max-Cut asks for a partition of the
vertices maximizing the weight of edges crossing the partition. A cut is
encoded by signs on vertices, and an edge contributes exactly when its signs
differ. For a maximization problem, an $\alpha$-approximation returns a
feasible solution worth at least $\alpha$ times the optimum. Unless stated
otherwise, the graph and its rational weights are given explicitly.
We write $\operatorname{MaxCut}(G)$ for the maximum crossing weight and,
when the edge weights sum to one, also denote this value by $\Val(G)$.
The decision problem of whether a weighted graph has a cut of at least a
given weight appears among Karp's original NP-complete
problems~\cite{Karp72}. Approximation asks how much of this optimum can
still be guaranteed in polynomial time.
Goemans and Williamson~\cite{GW95} gave a semidefinite programming algorithm with random-hyperplane rounding
for Max-Cut, with expected approximation ratio approaching
\begin{equation}\label{eq:gw-constant}
 \GW=\min_{-1\le \rho<1}
 \frac{2\arccos\rho}{\pi(1-\rho)}
 =0.878567\ldots.
\end{equation}
The expression compares two contributions of an edge whose endpoint
vectors have inner product $\rho$: its semidefinite objective value is
$(1-\rho)/2$, whereas a random hyperplane separates the vectors with
probability $\arccos(\rho)/\pi$. The minimum ratio over $\rho$ is the
uniform guarantee of this rounding rule.
Determining whether this ratio is optimal under ordinary NP-hardness has
been a central question in approximation algorithms. H\aa stad's
hardness theorem~\cite{Hastad01}, combined with the gadgets of
Trevisan, Sorkin, Sudan, and Williamson~\cite{TSSW00}, rules out ratios
greater than $16/17$.
Khot, Kindler, Mossel, and O'Donnell~\cite{KKMO07} established optimality of
$\GW$ assuming the Unique Games Conjecture~\cite{Khot02}; the analytic input
is the Majority Is Stablest theorem of Mossel, O'Donnell, and
Oleszkiewicz~\cite{MOO10}.

We resolve the unconditional optimality question for Max-Cut.

\begin{theorem}\label{thm:main}
For every fixed $\alpha$ with $\GW<\alpha\le1$, it is NP-hard to approximate
Max-Cut within a factor of $\alpha$. This holds for simple unweighted graphs.
\end{theorem}

The statement means that a polynomial-time algorithm which always returns a
cut of weight at least $\alpha$ times the maximum would imply
$\mathrm P=\mathrm{NP}$. We prove the following more precise gap theorem.
Write
\begin{equation}\label{eq:B}
 B(t)=\frac{2}{\pi}\arcsin t,\qquad 0<t<1.
\end{equation}

\begin{theorem}\label{thm:gap}
Fix a rational $t\in(0,1)$ and a constant
$0<\eps<(t-B(t))/4$. For graphs equipped with a rational probability
distribution on edges, it is NP-hard to distinguish
\begin{align*}
 \textnormal{YES:}\quad&\Val\ge\frac{1+t}{2}-\eps,\\
 \textnormal{NO:}\quad&\Val\le\frac{1+B(t)}{2}+\eps,
\end{align*}
where $\Val$ is the maximum cut weight.
\end{theorem}

\subsection*{Prior work and the direct route}

The reduction of proof verification to approximation gaps has its roots in
Feige, Goldwasser, Lov\'asz, Safra, and Szegedy~\cite{FGLSS96} and in the
PCP theorems of Arora--Safra and Arora, Lund, Motwani, Sudan, and
Szegedy~\cite{AroraSafra98,ALMSS98}. H\aa stad's Fourier analysis of
proof-verification tests~\cite{Hastad01} made many of these gaps sharp.
For Max-Cut, the Gaussian geometry of the semidefinite relaxation gives
a particularly precise target. Feige and Schechtman~\cite{FS02}
constructed integrality gaps approaching $\GW$ for this relaxation;
a limitation of that relaxation alone does not establish hardness for
arbitrary algorithms. The Unique Games reduction of
Khot, Kindler, Mossel, and O'Donnell~\cite{KKMO07} connected the same
geometric threshold to computational hardness.

The original analysis of Khot, Kindler, Mossel, and O'Donnell proposed
Majority Is Stablest as the analytic statement needed to obtain this
threshold. Mossel, O'Donnell, and Oleszkiewicz proved it through a
comparison between low-influence Boolean functions and Gaussian
functions~\cite{MOO10}, transferring Borell's halfspace noise-stability
bound~\cite{Borell85} to the Boolean setting. Their theorem leaves the
Unique Games Conjecture as the hypothesis in the classical hardness
reduction~\cite[Corollary~2.12]{MOO10}. We use Majority Is Stablest to
identify an influential coordinate of our test, then prove the additional
extraction needed to turn that coordinate into a source-game label.
The algorithmic benchmark above is the ideal rounding ratio;
finite-precision SDP optimization and random-hyperplane rounding give
expected ratio $\GW-\eta$ for any fixed $\eta>0$~\cite[Section~2]{GW95}.

The comparison between semidefinite approximation and hardness extends
beyond a single worst-case ratio. O'Donnell and Wu~\cite{OW08}
determined the full Max-Cut approximation curve under the Unique Games
Conjecture, relating the attainable cut value to the instance's optimum.
Raghavendra~\cite{Raghavendra08} established a general correspondence
between semidefinite relaxations and Unique-Games-based hardness for
finite constraint satisfaction problems. Our theorem concerns the
unconditional worst-case ratio for Max-Cut; the proof below establishes
the specific family of completeness and soundness gaps in
\Cref{thm:gap}.

In a projection game, each vertex of a bipartite graph receives a label
from a specified alphabet, and the edge weights form a probability
distribution. An edge accepts when its prescribed map sends the left label to the
right label. Its value is the maximum accepted edge weight. A two-to-one
constraint has exactly two preimages for every right label. Completeness
and soundness are, respectively, the promised lower bound on the YES-case
value and upper bound on the NO-case value.

Our starting point is the published 2-to-1 game construction with imperfect
completeness~\cite{DKKMS25}. The Grassmann-graph approach was introduced
by Khot, Minzer, and Safra~\cite{KMS25}. Its soundness was completed
through the connection between expansion and agreement established by
Barak, Kothari, and Steurer~\cite{BKS19} and the expansion theorem of
Khot, Minzer, and Safra~\cite{KMS23}. We need a precise affine form of the
construction: its projections are surjective affine maps between binary
spaces of dimensions differing by one, and its alphabet can be fixed for a
desired soundness before the completeness error is made arbitrarily small.
\Cref{prop:affine-games} states this input. Appendix~\ref{sec:affine-games}
derives it from the published construction, specifying both the
valid-subspace conditioning and the auxiliary candidate lists used in
the soundness proof.
Together with Majority Is Stablest, that construction supplies the
external foundation of the argument. The Unique Games Conjecture enters
only as historical context. The tree Fourier analysis, hashing bound,
and hidden-coordinate decoding are proved here.

\subsection*{The proof and its main ingredients}

As in long-code reductions~\cite{BGS98,Hastad01}, a label $a$ is represented
by the dictator $w\mapsto w(a)$, where $w$ ranges over Boolean words
indexed by labels. The graph has one such word domain for each tuple of
left vertices of the source game. Thus an arbitrary cut specifies a
Boolean function on each domain. To sample an edge, the test chooses two
left tuples through a common right tuple, generates two words on the
right label space, and pulls them back through the sampled affine
projections. It negates the second word. When the sampled constraints
are satisfied, dictator labels give a cut probability of one half of one
plus the correlation between the two words evaluated at a common label.

The word generator is a finite rooted tree. Every internal node has $s$
pairs of children and a random table from $\F^s$ to $\bits$; the leaves
carry affine forms on the label space. At a label $a$, these forms give
leaf signs, and the tables recursively combine the child signs into the
root sign $W(a)$. The gate rule is chosen so that each monomial of its
Walsh expansion keeps one child from every pair. A second exploration
keeps both children of one randomly selected pair at each reached node.
At any fixed depth, the two explorations have exactly one node in
common. The test couples the tables reached by this second exploration
at a random depth, using entrywise correlation $t$, and leaves all other
tables unchanged. Every monomial therefore acquires precisely one factor
of $t$. A small independent resampling of the affine rows gives the
smoothing needed for soundness and only a small completeness loss.

For an arbitrary cut, we first average its odd part over the source-game
edges. The resulting function is balanced in the selected table entries,
and a cut above the desired NO threshold makes its expected noise
stability exceed $B(t)$. Majority Is Stablest then detects an influential
entry. That entry is indexed by a gate input, not by a source-game label.
The remaining task is to connect its influence to labels that the two
endpoints of a source edge can choose consistently. Three steps address
this task.

\begin{enumerate}
\item \emph{Local simplification of Fourier support.}
The gate symmetries force the two child indices in each pair to have the
same sum. After row smoothing, averaging the depth of the selected tables
makes it inexpensive to retain indices which put this sum in one child of
each pair and zero in the other. This simplification is imposed only at
the influential gate. It is the structure that lets affine shifts act as
translations of that gate's table.

\item \emph{A vector-valued affine hashing lemma.}
For independent affine shifts, a diffuse distribution of Fourier mass
cannot create a large entry influence on average. The bound is uniform in
the dimension of the row space. Its proof uses two truncations and a
fourth-moment identity; it may be useful whenever translated coefficient
arrays must be decoded without an alphabet-size loss.

\item \emph{Decoding with a hidden coordinate.}
The mass concentration is obtained after fixing many row coordinates.
A left endpoint cannot usually construct the pullback of this context
without knowing its incident projection. We use a product of the game
and plant zero linear coefficients in one uniformly chosen factor of
the context and of the latent resampling rows. After that factor is
forgotten, the planted law is close to the original law. At the planted
factor, the context can be pulled back without knowing the projection.
Only the $2s$ free rows at the selected gate still have to be compared
across it, at a cost of $2^{2s}$ in Fourier mass.
\end{enumerate}

The resulting decoding loss is independent of the source alphabet.
This is essential: the alphabet dimension depends on the soundness we
request, so a dimension-dependent loss could make the soundness choice
circular. We first fix that loss and the source soundness, then the
alphabet and the number of product factors, and finally the source
completeness error. The quantifier order in
\Cref{prop:affine-games} permits exactly this sequence.

\Cref{sec:preliminaries} fixes the Fourier conventions.
\Cref{sec:tree,sec:reduction} define the code and the Max-Cut test.
\Cref{sec:tree-fourier,sec:hashing,sec:extraction} prove the analytic ingredients.
\Cref{sec:decoding,sec:completion} give the decoding and finish the proof.
Appendix~\ref{sec:affine-games} derives the precise starting game from the
published construction, and Appendix~\ref{sec:unweighted} gives a deterministic
conversion to simple unweighted graphs.
